% Lineare Funktionen · Ablesen · Prüfungs-Fokus MSA – bank/lineare-funktionen/e2.jsonl (zusammenbau v0.6, Prüfungs-Fokus)
\einheitenkopf[e2]{Einheit 2 · Lineare Funktion f(x) = m·x + n}
\zweigzeile{ab Kl. 8 · P10 ×4}
\setcounter{aufgabe}{0}

% Anlauf (ohne Prüfkennung)
\begin{aufgabe}{Ablesen – Anlauf}
\begin{teile}
\teil n? Lies den y-Achsenabschnitt n der Geraden g ab.

\begin{ksys}[xmin=-5,xmax=5,ymin=-5,ymax=5,ablesen] \gerade{1}{3}{g} \end{ksys}
n = \leerfeld
\teil m? Lies die Steigung m der Geraden g mit einem Steigungsdreieck ab.

\begin{ksys}[xmin=-5,xmax=5,ymin=-5,ymax=5,ablesen] \gerade{-2}{1}{g} \end{ksys}
m = \leerfeld
\end{teile}
\end{aufgabe}

\begin{aufgabe}{Ablesen – Prüfungsaufgaben}
\begin{teile}
\teil Welche Gleichung gehört zu welcher Geraden? Ordne den Geraden f, g und h je eine der Gleichungen zu: $y = -3x + 3$; $y = -\frac{1}{3}x + 3$; $y = -3$; $y = 2x + 3$; $y = -3x$; $y = 2x - 3$. \hfill (P19) \\

\begin{ksys}[xmin=-5,xmax=5,ymin=-5,ymax=5,ablesen] \gerade{-3}{3}{f} \gerade{2}{3}{g} \gerade{0}{-3}{h} \end{ksys}
f: \leerfeld \quad g: \leerfeld \quad h: \leerfeld
\teil Welche Gleichung gehört zu welcher Geraden? Ordne den Geraden f, g und h je eine der Gleichungen zu: $y = -4x + 4$; $y = -\frac{1}{4}x + 4$; $y = -1$; $y = 2x + 4$; $y = -4x$; $y = 2x - 4$. \hfill (P19) \\

\begin{ksys}[xmin=-5,xmax=5,ymin=-5,ymax=5,ablesen] \gerade{-4}{4}{f} \gerade{2}{4}{g} \gerade{0}{-1}{h} \end{ksys}
f: \leerfeld \quad g: \leerfeld \quad h: \leerfeld
\teil Welche Gerade hat die Eigenschaft? Gib zu jeder Eigenschaft die passende Gerade an: Anstieg $-3$; parallel zur x-Achse. \hfill (P19) \\

\begin{ksys}[xmin=-5,xmax=5,ymin=-5,ymax=5,ablesen] \gerade{-3}{1}{f} \gerade{3}{1}{g} \gerade{0}{3}{h} \end{ksys}
Anstieg: \leerfeld \quad parallel zur x-Achse: \leerfeld
\teil Welche Gerade hat die Eigenschaft? Gib zu jeder Eigenschaft die passende Gerade an: Anstieg $-2$; parallel zur x-Achse. \hfill (P19) \\

\begin{ksys}[xmin=-5,xmax=5,ymin=-5,ymax=5,ablesen] \gerade{2}{-1}{f} \gerade{-2}{-1}{g} \gerade{0}{-3}{h} \end{ksys}
Anstieg: \leerfeld \quad parallel zur x-Achse: \leerfeld
\teil Welche Gerade fällt? Kreuze an. \hfill (P16) \\ \kreuz{f} \\ \kreuz{g}

\begin{ksys}[xmin=-5,xmax=5,ymin=-5,ymax=5,ablesen] \gerade{1.5}{-1}{f} \gerade{-0.5}{2}{g} \end{ksys}
\teil Welche Gerade fällt? Kreuze an. \hfill (P16) \\ \kreuz{f} \\ \kreuz{g}

\begin{ksys}[xmin=-5,xmax=5,ymin=-5,ymax=5,ablesen] \gerade{-1}{3}{f} \gerade{2}{-3}{g} \end{ksys}
\teil Schnittpunkt mit der y-Achse: Markiere ihn an der Geraden $y = 1{,}5x - 3$. \hfill (P24)

\begin{ksys}[xmin=-5,xmax=5,ymin=-5,ymax=5,ablesen] \gerade{1.5}{-3}{g} \end{ksys}
\teil Schnittpunkt mit der y-Achse: Markiere ihn an der Geraden $y = -1{,}5x + 3$. \hfill (P24)

\begin{ksys}[xmin=-5,xmax=5,ymin=-5,ymax=5,ablesen] \gerade{-1.5}{3}{g} \end{ksys}
\teil Welcher Graph gehört zu einer linearen Funktion? Kreuze an. \hfill (P25) \\ \kreuz{A} \\ \kreuz{B} \\ \kreuz{C} \\ \kreuz{D}

\begin{ksys}[xmin=-5,xmax=5,ymin=-5,ymax=5,ablesen] \gerade{-0.5}{3}{A} \funktionab{2/\x}{B}{0.4}{5} \parabel{1}{2}{-4}{C} \funktion{abs(\x+2)-4}{D} \end{ksys}
\teil Welcher Graph gehört zu einer linearen Funktion? Kreuze an. \hfill (P25) \\ \kreuz{A} \\ \kreuz{B} \\ \kreuz{C} \\ \kreuz{D}

\begin{ksys}[xmin=-5,xmax=5,ymin=-5,ymax=5,ablesen] \funktion{abs(\x-2)+1}{A} \parabel{-1}{-2}{4}{B} \gerade{1.5}{-3}{C} \funktionab{-3/\x}{D}{0.6}{5} \end{ksys}
\end{teile}
\end{aufgabe}

\medskip
\uebersichtskasten{Lineare Funktion: f(x) = m · x + n \\ m = Steigung: 1 nach rechts, m nach oben (m < 0: nach unten). \quad f(x) = 2x + 1: 1 nach rechts, 2 nach oben \\ n = y-Achsenabschnitt: die Gerade schneidet die y-Achse bei (0 | n). \quad f(x) = 2x + 1: bei (0 | 1)}
